\scanpage{4}
\mspage{2}

1.3. Regardons de même ce qui se passe quand on fixe $\underline F$ et fait varier R,S. Tout d'abord, si $S\to R'\to R$, alors tout cristal sur S relativement à R en définit un \deleted{sur} relativement à R', par simple restriction. En particulier, un cristal absolu définit un cristal relatif, pour tout préschéma de base R de S.

Fixons maintenant R,S et soit $S'\to S$ un morphisme. On définit alors de façon naturelle un foncteur ``image inverse'' allant des cristaux de type $\underline F$ sur S vers les cristaux de même type sur S' (tout relatif à R). Pour s'en assurer, \unsure{il|id} suffit de définir un foncteur $C'\longrightarrow C$ (ou $C'\deleted{\unsure{a}}$ est défini en termes de S' comme C en termes de S), compatible avec les foncteurs ``oubli''. Or si $S'\longrightarrow T'$ est un objet de C', il y a une construction évidente de somme amalgamée dans la catégorie des préschémas, qui donne un $T=S\amalg_{S'}T'$ et un morphisme $S\longrightarrow T$, qui fait de T un objet de C, ce qui définit le foncteur cherché.

On peut donc dire que pour un S variable sur R, les cristaux de type $\underline F$ sur S forment une \uline{catégorie fibrée} sur $(\mathrm{Sch})_{/R}$, grâce à la notion d'image inverse précédente.

Les deux variances (en R, et en S) sont compatibles dans un sens évident, et peuvent être regardées comme provenant d'une variance en (R,S) directement.

1.4. On a un foncteur naturel et évident qui va des \unsure{criataux|cristaux} de modules (disons) sur S (rel à R) dans des Modules quasi-cohérents sur S: c'est le foncteur ``valeur en S''. On fera attention que ce foncteur n'est en général pas même fidèle (cf exemple \inserted{1.5.} plus bas). Il est donc un peu dangereux de vouloir considérer un cristal de modules sur S comme étant un Module quasi-cohérent $\underline M$ sur S, \uline{muni} d'une structure \unsure{spplémentaire|supplémentaire}, sa ``structure cristalline''. Dans certains cas cependant \inserted{(cf 1.8.)}, le foncteur cristaux de modules $\rightsquigarrow$Modules est fidèle et le point de vue précédent devient plus raisonnable,

1.5. \uline{Exemple} 1: \uline{relations avec les vecteurs de Witt}. Supposons que S soit le spectre d'un corps \uline{parfait} \deleted{\unsure{du}} k, et $R=\operatorname{Spec}(\underline Z)$. Soit W l'anneau de Cohen de corps résiduel k: si $\operatorname{car}k>0$, c'est l'anneau des vecteurs de Witt défini par k, si $\operatorname{car}.k=0$, c'est k lui-même\inserted{; dans ce dernier cas, supposons k \uline{algébrique} sur Q.} Alors il est bien connu que la catégorie C de 1.0.
\margin{3}
